\begin{table}[!htbp]
\caption{Сопоставление с ближайшими работами по компиляции и проверке. Настоящая работа добавляет подтверждения для конкретной реализации, а не новую общую теорию верификации.}
\label{tab:related-work}
\centering\scriptsize\setlength{\tabcolsep}{3pt}
\begin{tabular}{p{2.2cm}p{3.25cm}p{3.1cm}p{2.6cm}p{4.3cm}}
\toprule
Работа & Исходное представление и табуляция & Гарантия / качество & Модель и задача & Аппаратная область, размещение, энергия\\
\midrule
LUT-KAN~\cite{kuznetsov2026lutkan} & Сплайны или полные функции рёбер PyKAN; сегментные int8/uint8-таблицы с интерполяцией & Эмпирические ошибки и сохранение предсказаний; явные правила границ & Контролируемые слои KAN и DoS CICIDS2017 & Настольный CPU; сравнение сплайнов/LUT при одинаковом NumPy- или Numba-бэкенде.\\
LUT-compiled DoS~\cite{kuznetsov2026lutdos} & Многослойная B-сплайновая KAN; квантованные отсчёты с интерполяцией & Эмпирический компромисс точности, памяти и задержки & Бинарное обнаружение DoS в CICIDS2017 & Настольный CPU (NumPy/Numba); энергия обсуждается через задержку как прокси, без прямого измерения мощности.\\
KANEL\'{E}~\cite{Hoang2025KANELE} & Полные обученные одномерные функции в LUT & Аппаратное отображение и эмпирическое качество & MNIST, классификация струй, табличные и аудиозадачи & FPGA с квантованием и прореживанием; здесь интерполяция отсчётов и согласованность с integer-эталоном на измеренных MCU.\\
QuantKAN / KANtize~\cite{fuad2025quantkan,errabii2026kantize} & Низкоразрядное квантование; табуляция B-сплайнов и полных функций & Качество после квантования и стоимость & Квантованные KAN & Контекст развёртывания; здесь полный перебор знаковой ошибки конкретного вычислителя.\\
QEBVerif~\cite{zhang2023qebverif} & Сеть и её квантованный аналог & Границы ошибки квантования & Квантованные сети & Контекст формальной проверки; здесь используется аддитивная структура конечной области.\\
Эта работа & Зафиксированные integer-коэффициенты; отсчёты LUT с интерполяцией & Точные интервалы рёбер; достаточные сертификаты пороговых решений & Одна аддитивная оценка TON\_IoT & Mega 2560 и ESP32-C3; контролируемое размещение кода/параметров; отдельный пилот энергии всей платы через USB.\\
\bottomrule
\end{tabular}
\end{table}
